\chapter{Safety, Correctness, and Concurrency}
\label{chap:ch5}

\par Modern banking backends must stay correct under load, not only crash-free. Gentlix Bank
implements the same product in C and Rust, but the two codebases prove safety in different
places: C through conventions and review, Rust through types and the borrow checker.
This chapter examines three topics that sit on top of the design patterns from
Chapter~\ref{chap:ch4}: \textbf{type-driven invariants} in the domain layer,
\textbf{error propagation} through services, and \textbf{concurrency} in the analytics
pipeline. Database schema alignment and HTTP error envelopes belong to
Chapter~\ref{sec:ch6sec1} and Section~\ref{subsec:ch6sec2-5}; they are referenced here
only where a service decision depends on them.

\section{Type-Driven Invariants}
\label{sec:ch5sec1}

Financial software can return HTTP success with wrong numbers. Gentlix therefore treats
\emph{type-driven invariants} as first-class rules: what a value means, which combinations
are illegal, and which constructor must be used before data enters the service layer.
Blackburn and colleagues~\cite{Blackburn2025} separate spatial and temporal memory bugs
from domain mistakes; this section focuses on the domain layer in \texttt{domain/} and
\texttt{service/}---money, identities, and transaction shape.

In C, many invariants are enforced by helper functions and repeated \texttt{if} checks.
The compiler sees \texttt{int64\_t} and \texttt{const char *} whether or not a value was
validated. In Rust, invalid states are often \emph{unrepresentable}: a value simply
cannot exist unless it passed a validating constructor. Palmieri~\cite{Palmieri2022}
describes this as the central idea of type-driven development---``we made it impossible
for an instance of \texttt{SubscriberName} to violate those constraints [\ldots] it will
not compile.'' Gentlix applies the same philosophy to \texttt{Email}, \texttt{IBAN}, and
account identifiers.

\medskip\noindent\textbf{Money, value objects, and identifiers.}\quad
Floating-point must never represent customer balances. Gentlix stores money as
\texttt{int64\_t} cents end-to-end. ATM-style deposits accept whole currency units in JSON
but convert before domain logic runs. Rust rejects overflow explicitly with
\texttt{checked\_mul}; C guards with an explicit bound before multiplying:

\begin{lstlisting}[language=Rust, caption={transaction.rs (service) --- Unit-to-cent
conversion uses \texttt{checked\_mul}; wraparound becomes a validation error.}]
let value_cents = amount_units
    .checked_mul(100)
    .ok_or_else(|| ServiceError::Validation("amount too large".into()))?;
\end{lstlisting}

\begin{lstlisting}[language=C, caption={transaction\_service.c --- Same rule: reject
amounts that would overflow when converted to cents.}]
if (amount_units > INT64_MAX / 100) {
    if (err) *err = service_error_validation("amount too large");
    return false;
}
cmd.value_cents = amount_units * 100;
\end{lstlisting}

Drysdale~\cite{Drysdale2024}
recommends making monetary arithmetic explicit rather than assuming inputs stay small.

Value objects wrap raw strings after validation. In Rust, \texttt{Email} is built only
through \texttt{FromStr}:

\begin{lstlisting}[language=Rust, caption={email.rs --- An \texttt{Email} cannot exist
without passing validation.}]
impl FromStr for Email {
    type Err = DomainError;
    fn from_str(s: &str) -> Result<Self, Self::Err> {
        let s = s.trim();
        if s.is_empty() {
            return Err(DomainError::validation("Email must not be empty"));
        }
        if s.split('@').collect::<Vec<_>>().len() != 2 {
            return Err(DomainError::validation(
                "Email must contain exactly one '@'"));
        }
        Ok(Email(s.to_string()))
    }
}
\end{lstlisting}

In the login flow, \texttt{user.rs} parses once and shadows the raw string
(Section~\ref{sec:ch3sec3}). C validates through \texttt{email\_try\_create} and
\texttt{iban\_validate}, but a validated \texttt{Email} struct still stores a plain
\texttt{char[]} that looks like any other string elsewhere in the file:

\begin{lstlisting}[language=C, caption={email.c --- Validation is explicit; the return
type does not carry the proof to distant call sites.}]
bool email_try_create(const char *raw, Email *out, DomainError *err) {
    /* trim whitespace, reject empty and overlong input */
    char *at = strchr(buffer, '@');
    if (!at || strchr(at + 1, '@')) {
        if (err) *err = domain_error_validation(
            "Email must contain exactly one '@'");
        return false;
    }
    strcpy(out->value, buffer);
    return true;
}
\end{lstlisting}

\texttt{IBAN} follows the same value-object pattern. Rust validates length and country
code in \texttt{FromStr}; C normalises whitespace and checks the same invariants in
\texttt{iban\_try\_create}:

\begin{lstlisting}[language=Rust, caption={iban.rs --- Invalid IBAN strings cannot
become an \texttt{IBAN} value.}]
impl FromStr for IBAN {
    fn from_str(s: &str) -> Result<Self, Self::Err> {
        let iban: String = s.trim().chars()
            .filter(|c| !c.is_whitespace()).collect();
        if iban.len() != 24 {
            return Err(DomainError::validation("Invalid IBAN length"));
        }
        if !iban.starts_with("RO") {
            return Err(DomainError::validation("Invalid country code"));
        }
        Ok(IBAN(iban))
    }
}
\end{lstlisting}

\begin{lstlisting}[language=C, caption={iban.c --- Same constraints; proof lives in
the helper, not the C type system.}]
bool iban_try_create(const char *raw, IBAN *out, DomainError *err) {
    /* strip whitespace, uppercase, require len == 24 and RO prefix */
    if (strlen(buffer) != 24) { /* ... */ return false; }
    if (buffer[0] != 'R' || buffer[1] != 'O') { /* ... */ return false; }
    strcpy(out->value, buffer);
    return true;
}
\end{lstlisting}

Drysdale~\cite{Drysdale2024} recommends pushing validation to the edges of a system
so invalid data never reaches core logic. Gentlix does that in both languages; Rust
simply makes the edge harder to bypass.

Swapping \texttt{UserId} and \texttt{AccountId} is an access-control bug. C names them
as separate structs in \texttt{ids.h}, yet both wrap \texttt{int64\_t value}:

\begin{lstlisting}[language=C, caption={ids.h --- Intent is documented; the compiler
does not prevent mixing identifiers.}]
typedef struct { int64_t value; } UserId;
typedef struct { int64_t value; } AccountId;
\end{lstlisting}

Rust newtypes are incompatible at compile time:

\begin{lstlisting}[language=Rust, caption={ids.rs --- \texttt{UserId} and
\texttt{AccountId} are distinct types.}]
pub struct UserId(pub i64);
pub struct AccountId(pub i64);
\end{lstlisting}

\texttt{ensure\_account\_owned\_by(user\_id, account\_id)} in Rust cannot compile with
reversed arguments. Both backends load the account and compare \texttt{user\_id} before
any debit:

\begin{lstlisting}[language=C, caption={transaction\_service.c --- Ownership check
before \texttt{record\_transfer\_for\_user} delegates to the core recorder.}]
static bool ensure_account_owned_by(
    TransactionService *svc, UserId user_id,
    AccountId account_id, ServiceError *err)
{
    Account *account = NULL;
    if (!svc->account_repo.vtable->get_by_id(
            svc->account_repo.ctx, account_id, &account, &rerr))
        return false;
    bool owned = (account_user_id(account).value == user_id.value);
    account_free(account);
    if (!owned) {
        if (err) *err = service_error_not_found("account");
        return false;
    }
    return true;
}
\end{lstlisting}

\begin{lstlisting}[language=Rust, caption={transaction.rs (service) --- Same rule;
wrong owner becomes \texttt{NotFound}, not a debit attempt.}]
async fn ensure_account_owned_by(
    &self, user_id: UserId, account_id: AccountId,
) -> ServiceResult<()> {
    let account = self.account_repo.get_by_id(account_id).await?;
    if account.user_id() != user_id {
        return Err(ServiceError::not_found("account"));
    }
    Ok(())
}
\end{lstlisting}

The C equivalent at call sites relies on review across dozens of paths.

\medskip\noindent\textbf{Transaction and account rules.}\quad
A transaction with the same source and destination, or with a negative amount, must not
reach persistence. Rust centralises these rules in \texttt{Transaction::build}:

\begin{lstlisting}[language=Rust, caption={transaction.rs (domain) --- Illegal states
are rejected before a \texttt{Transaction} value exists.}]
if from_account_id == to_account_id {
    return Err(DomainError::validation(
        "From Account ID must be different from To Account ID"));
}
if value_cents < 0 {
    return Err(DomainError::validation("Value must be >= 0"));
}
\end{lstlisting}

C's \texttt{transaction\_create} enforces the same constraints and returns \texttt{NULL}
with a \texttt{DomainError} out-parameter:

\begin{lstlisting}[language=C, caption={transaction.c --- Domain rules live in one
function; callers must check the return pointer.}]
Transaction *transaction_create(
    AccountId from, AccountId to, TransactionType type,
    int64_t value_cents, const char *desc, DomainError *err)
{
    if (from.value == to.value) {
        if (err) domain_error_set(err, DOMAIN_ERROR_VALIDATION,
            "From Account ID must be different from To Account ID");
        return NULL;
    }
    if (value_cents < 0) { /* ... */ return NULL; }
    /* allocate and populate Transaction */
}
\end{lstlisting}

The service layer in both languages adds \texttt{value\_cents <= 0} checks again before
calling the domain constructor:

\begin{lstlisting}[language=C, caption={transaction\_service.c --- Service rejects
non-positive amounts before \texttt{transaction\_create}.}]
if (cmd->value_cents <= 0) {
    if (err) *err = service_error_validation(
        "transaction value must be strictly positive");
    return false;
}
\end{lstlisting}

\begin{lstlisting}[language=Rust, caption={transaction.rs (service) --- Same guard
at the top of \texttt{record\_transaction}.}]
if cmd.value_cents <= 0 {
    return Err(ServiceError::Validation(
        "transaction value must be strictly positive".into()));
}
\end{lstlisting}

Defence in depth by convention in C; redundant but harmless when Rust already narrows
the input through earlier validation.

A row that has not yet been inserted has no database id. Rust makes that state explicit
through \texttt{Option<TransactionId>}: \texttt{Transaction::create} passes \texttt{None};
\texttt{Transaction::rehydrate} passes \texttt{Some(id)} after a PostgreSQL read:

\begin{lstlisting}[language=Rust, caption={transaction.rs (domain) --- \texttt{None}
means ``not persisted yet''; \texttt{Some} means reloaded from the database.}]
pub struct Transaction {
    id: Option<TransactionId>,
    from_account_id: AccountId,
    to_account_id: AccountId,
    /* ... */
}

pub fn create(/* ... */) -> Result<Self, DomainError> {
    Self::build(None, from_account_id, /* ... */)
}

pub fn id(&self) -> Option<TransactionId> {
    self.id
}
\end{lstlisting}

C stores the same distinction as a \texttt{bool has\_id} flag beside a
\texttt{TransactionId} field. \texttt{transaction\_create} sets \texttt{has\_id = false};
the repository refuses insert when the flag is already set:

\begin{lstlisting}[language=C, caption={transaction.c / transaction\_repo.c ---
C uses a flag; Rust uses \texttt{Option}. Both reject double-insert.}]
struct Transaction {
    bool has_id;
    TransactionId id;
    /* ... */
};

// transaction_repo.c --- guard before INSERT
if (transaction_has_id(tx)) {
    if (err) *err = repo_error_from_domain(&derr);
    return false;   /* already has an id */
}
\end{lstlisting}

Nothing in C stops a caller from reading \texttt{transaction\_id(tx)} when
\texttt{has\_id} is false---the value may be zero or stale. Rust forces callers through
\texttt{id()} and pattern-match on \texttt{Option}.

Opening a second account with the same \texttt{(account\_type, currency)} pair must fail
with a conflict error. Both backends query the repository before insert:

\begin{lstlisting}[language=C, caption={account\_service.c --- Duplicate
(type, currency) rejected before domain \texttt{Account} is built.}]
if (type_currency_exists) {
    snprintf(msg, sizeof msg,
             "you already have a %s %s account",
             account_type_as_str(cmd->account_type),
             currency_as_str(cmd->currency));
    if (err) *err = service_error_conflict("account", msg);
    return false;
}
\end{lstlisting}

\begin{lstlisting}[language=Rust, caption={account.rs (service) --- Same rule via
\texttt{ServiceError::Conflict}.}]
if self.repo.exists_by_type_and_currency(
    user_id, account_type, currency).await? {
    return Err(ServiceError::conflict("account", format!(
        "you already have a {} {} account",
        account_type.as_str(), currency.as_str())));
}
\end{lstlisting}

When a user submits \texttt{POST /api/transactions/transfer}, the service rejects
non-positive amounts, verifies account ownership, checks overdraft before debit, and only
then asks the repository to persist. The \texttt{record\_transfer\_for\_user} wrappers
call \texttt{ensure\_account\_owned\_by} on both accounts first:

\begin{lstlisting}[language=C, caption={transaction\_service.c --- Transfer route
checks ownership before recording.}]
bool transaction_service_record_transfer_for_user(
    TransactionService *svc, UserId user_id,
    AccountId from_id, AccountId to_id,
    int64_t value_cents, TransactionId *out_id, ServiceError *err)
{
    if (!ensure_account_owned_by(svc, user_id, from_id, err)) return false;
    if (!ensure_account_owned_by(svc, user_id, to_id, err))   return false;
    return record_transfer(svc, from_id, to_id, value_cents, out_id, err);
}
\end{lstlisting}

\begin{lstlisting}[language=Rust, caption={transaction.rs (service) --- Same wrapper
shape.}]
pub async fn record_transfer_for_user(
    &self, user_id: UserId, from_id: AccountId,
    to_id: AccountId, value_cents: i64,
) -> ServiceResult<TransactionId> {
    self.ensure_account_owned_by(user_id, from_id).await?;
    self.ensure_account_owned_by(user_id, to_id).await?;
    self.record_transfer(from_id, to_id, value_cents).await
}
\end{lstlisting}

Both backends compare the source account balance
against the requested amount \emph{before} calling \texttt{account\_debit} or
\texttt{Account::debit}:

\begin{lstlisting}[language=C, caption={transaction\_service.c --- Overdraft rejected
in the service layer before any balance mutation.}]
if (account_balance_cents(from_acc) < cmd->value_cents) {
    account_free(from_acc);
    if (err) *err = service_error_validation("insufficient funds");
    return false;
}
if (!account_debit(from_acc, cmd->value_cents, &derr)) { /* ... */ }
\end{lstlisting}

\begin{lstlisting}[language=Rust, caption={transaction.rs (service) --- Same guard on
transfer and withdrawal paths.}]
if from_acc.balance_cents() < cmd.value_cents {
    return Err(ServiceError::Validation("insufficient funds".into()));
}
from_acc.debit(cmd.value_cents)?;
\end{lstlisting}

The domain layer repeats the rule if a future caller skips the service check:

\begin{lstlisting}[language=C, caption={account.c --- \texttt{account\_debit} is the
last line of defence against overdraft.}]
if (a->balance_cents < amount_cents) {
    if (err) *err = domain_error_validation("Insufficient funds");
    return false;
}
a->balance_cents -= amount_cents;
\end{lstlisting}

\begin{lstlisting}[language=Rust, caption={account.rs (domain) --- Same invariant
inside \texttt{Account::debit}.}]
if self.balance_cents < amount_cents {
    return Err(DomainError::validation("Insufficient funds"));
}
self.balance_cents -= amount_cents;
\end{lstlisting}

Ledger atomicity at the database boundary is implemented in the repository layer
(Section~\ref{subsec:ch6sec1-1}); the service assumes debit, credit, and insert commit
together or roll back entirely.

\section{Error Propagation}
\label{sec:ch5sec2}

When withdrawal exceeds balance, Gentlix must surface a structured failure the React app
can display. Errors originate in the domain and repository layers and convert upward into
\texttt{ServiceError}. Mapping those enums to HTTP status codes happens only in the HTTP
layer (Section~\ref{subsec:ch6sec2-5}); this section stays inside \texttt{service/}.

Schreiner~\cite{Schreiner1993} describes the idiomatic C pattern: return a status flag
and optionally write into a caller-provided error struct---``messages to objects'' paired
with out-parameters. Gentlix follows that model in \texttt{service\_error.c} and every
service entry point:

\begin{lstlisting}[language=C, caption={transaction\_service.c --- The C service
returns \texttt{bool} and optionally fills \texttt{ServiceError *err}.}]
bool transaction_service_record_transaction(
    TransactionService             *svc,
    const RecordTransactionCommand *cmd,
    TransactionId                  *out_id,
    ServiceError                   *err)
\end{lstlisting}

The recording pipeline can fail at every step: domain creation, repository
\texttt{begin}, balance application, \texttt{insert}, or \texttt{commit}. C opens the
SQL transaction before touching balances:

\begin{lstlisting}[language=C, caption={transaction\_service.c --- \texttt{begin}
pairs balance updates and insert in one DB transaction.}]
if (!svc->tx_repo.vtable->begin(svc->tx_repo.ctx, &rerr)) {
    transaction_free(tx);
    if (err) *err = service_error_from_repo(&rerr);
    return false;
}
if (!record_transaction_apply_balances(svc, cmd, err)) {
    svc->tx_repo.vtable->rollback(svc->tx_repo.ctx);
    /* ... */
}
\end{lstlisting}

Rust delegates the equivalent boundary to \texttt{record\_atomic} (below). On failure,
C pairs resources on each branch---\texttt{transaction\_free} and \texttt{rollback}:

\noindent\begin{minipage}{\linewidth}
\begin{lstlisting}[language=C, caption={transaction\_service.c --- Error paths free
the in-memory row and roll back the SQL transaction.}]
if (!record_transaction_apply_balances(svc, cmd, err)) {
    svc->tx_repo.vtable->rollback(svc->tx_repo.ctx);
    transaction_free(tx);
    return false;
}
bool ok = svc->tx_repo.vtable->insert(
    svc->tx_repo.ctx, tx, &new_id, &rerr);
transaction_free(tx);
if (!ok) {
    svc->tx_repo.vtable->rollback(svc->tx_repo.ctx);
    if (err) *err = service_error_from_repo(&rerr);
    return false;
}
\end{lstlisting}\end{minipage}

Gatchell~\cite{Gatchell2003} stresses that missed cleanup branches are where production
C systems silently leak or corrupt state. The deeper weakness is optional checking: nothing
stops a caller from ignoring the \texttt{bool} return or passing \texttt{NULL} for
\texttt{err}. Domain and repository failures convert through explicit helpers in C:

\begin{lstlisting}[language=C, caption={service\_error.c / call sites --- Manual
enum lift on every branch.}]
ServiceError service_error_from_domain(const DomainError *derr);
ServiceError service_error_from_repo(const RepoError *rerr);

if (!tx) {
    if (err) *err = service_error_from_domain(&derr);
    return false;
}
\end{lstlisting}

Rust collapses the pattern into \texttt{ServiceResult<T>}. The \texttt{thiserror} crate
lifts \texttt{DomainError} and \texttt{RepoError} automatically:

\begin{lstlisting}[language=Rust, caption={errors.rs --- \texttt{\#[from]} replaces
manual conversion helpers.}]
#[derive(Debug, Error)]
pub enum ServiceError {
    #[error("validation error: {0}")]
    Validation(String),
    #[error(transparent)]
    Domain(#[from] DomainError),
    #[error(transparent)]
    Repo(#[from] RepoError),
    #[error("forbidden")]
    Forbidden,
    #[error("concurrency error: {0}")]
    Concurrency(String),
}
\end{lstlisting}

The same transfer path reads linearly with \texttt{?}:

\noindent\begin{minipage}{\linewidth}
\begin{lstlisting}[language=Rust, caption={transaction.rs (service) ---
\texttt{record\_transaction} propagates failures with \texttt{?}.}]
pub async fn record_transaction(
    &self, cmd: RecordTransactionCommand,
) -> ServiceResult<TransactionId> {
    if cmd.value_cents <= 0 {
        return Err(ServiceError::Validation(
            "transaction value must be strictly positive".into()));
    }
    let tx = Transaction::create(
        cmd.from_account_id, cmd.to_account_id,
        cmd.transaction_type, cmd.value_cents,
        cmd.description,
    )?;
    /* ... apply balances, insert via repository ... */
    Ok(id)
}
\end{lstlisting}\end{minipage}

Drysdale~\cite{Drysdale2024} notes that \texttt{Result} keeps the success path readable
compared with nested \texttt{if (!ok)} blocks. The C recording path continues after
balance updates with \texttt{commit}; any failure rolls back:

\begin{lstlisting}[language=C, caption={transaction\_service.c --- Commit only after
insert succeeds; otherwise rollback.}]
if (!ok) {
    svc->tx_repo.vtable->rollback(svc->tx_repo.ctx);
    if (err) *err = service_error_from_repo(&rerr);
    return false;
}
if (!svc->tx_repo.vtable->commit(svc->tx_repo.ctx, &rerr)) {
    svc->tx_repo.vtable->rollback(svc->tx_repo.ctx);
    if (err) *err = service_error_from_repo(&rerr);
    return false;
}
*out_id = new_id;
return true;
\end{lstlisting}

Rust routes persistence through \texttt{record\_atomic} on the repository trait---one
await that applies balance updates and insert inside a SQL transaction (details in
Section~\ref{subsec:ch6sec1-1}):

\begin{lstlisting}[language=Rust, caption={transaction.rs (service) --- Trait method
that wraps balance \texttt{UPDATE}s and row \texttt{INSERT} in one SQL transaction.}]
async fn record_atomic(
    &self,
    updated_accounts: &[Account],
    tx: &Transaction,
) -> Result<TransactionId, RepoError>;

// call site inside record_transaction:
let id = self.tx_repo.record_atomic(&updated_accounts, &tx).await?;
\end{lstlisting}

The service method ends with a single \texttt{Ok(id)}.

Login applies the same policy in both languages: invalid email format, missing user, and
wrong password all map to the same message so the API does not reveal whether an address
exists:

\begin{lstlisting}[language=C, caption={user\_service.c --- Uniform
\texttt{invalid credentials} on every failure path.}]
if (!email_try_create(cmd->email, &email, &derr)) {
    if (err) *err = service_error_validation("invalid credentials");
    return false;
}
if (!repo_ok) {
    if (err) *err = service_error_validation("invalid credentials");
    return false;
}
if (rc != ARGON2_OK) {
    if (err) *err = service_error_validation("invalid credentials");
    return false;
}
\end{lstlisting}

\begin{lstlisting}[language=Rust, caption={user.rs (service) --- Same three paths,
same string.}]
let email: Email = cmd.email.parse()
    .map_err(|_| ServiceError::Validation("invalid credentials".into()))?;
let user = self.user_repo.get_by_email(&email).await.map_err(|e| match e {
    RepoError::NotFound(_) =>
        ServiceError::Validation("invalid credentials".into()),
    other => ServiceError::from(other),
})?;
if !valid {
    return Err(ServiceError::Validation("invalid credentials".into()));
}
\end{lstlisting}

Palmieri~\cite{Palmieri2022} treats unhandled service failures as production incidents
waiting to happen. In C the risk is \textbf{silent success}; in Rust unused
\texttt{Result} values trigger warnings and \texttt{?} makes propagation the default.
Parallel analytics adds \texttt{ServiceError::Concurrency} when worker spawn or join
fails---mechanics in Section~\ref{sec:ch5sec3}.

\section{Concurrency Patterns}
\label{sec:ch5sec3}

Most Gentlix requests are sequential: one transfer, one account lookup. The heavy
exception is \texttt{compute\_user\_statistics}, which loads transactions for every
owned account, deduplicates by id, splits work across threads, and merges totals.
Statement replay and Argon2 login also consume CPU time, but they use a different
pattern---blocking the connection thread in C or \texttt{spawn\_blocking} in Rust
(Section~\ref{subsec:ch6sec2-4}). This section explains the three mechanisms the
analytics pipeline combines: \textbf{mutexes} for shared maps, \textbf{atomics} for
global counters, and \textbf{thread lifetimes} for safe parallel slices.

Both backends follow the same high-level design documented by Bos~\cite{Bos2023}:
keep critical sections small, prefer local accumulation, merge under a lock once per
worker. The difference is how much the language enforces.

\subsection{Mutexes and shared maps}
\label{subsec:ch5sec3mutex}

A mutex ensures only one thread mutates shared data at a time. In analytics, per-type and
per-day totals must not be updated concurrently without exclusion. Gentlix stores shared
state in \texttt{StatsShared}:

\begin{lstlisting}[language=C, caption={transaction\_service.c --- Atomics hold
global counters; mutexes protect richer maps.}]
typedef struct StatsShared {
    _Atomic int64_t total_incoming;
    _Atomic int64_t total_outgoing;
    _Atomic int64_t total_volume;
    _Atomic int64_t transaction_count;

    pthread_mutex_t per_type_mutex;
    PerTypeTotals   per_type;

    pthread_mutex_t per_day_mutex;
    DayTotalMap     per_day;
} StatsShared;
\end{lstlisting}

Each \texttt{stats\_worker\_fn} first accumulates into \textbf{thread-local} structures
without locking, then merges once:

\begin{lstlisting}[language=C, caption={transaction\_service.c --- Locks are held
only during the merge phase, not in the hot loop.}]
pthread_mutex_lock(&s->per_type_mutex);
for (size_t c = 0; c < TX_SERVICE_TYPE_COUNT; ++c) {
    if (local_type.present[c]) {
        s->per_type.totals[c] += local_type.totals[c];
        s->per_type.present[c]  = true;
    }
}
pthread_mutex_unlock(&s->per_type_mutex);
\end{lstlisting}

In C, the mutex and the data it protects live in separate fields---nothing forces a lock
before access except discipline. \texttt{pthread\_mutex\_init} and
\texttt{pthread\_mutex\_destroy} must run on every entry point. Rust wraps maps in
\texttt{Arc<Mutex<HashMap<TransactionType, i64>>>} and
\texttt{Arc<Mutex<BTreeMap<NaiveDate, i64>>>}; mutation without \texttt{lock()} does not
compile:

\begin{lstlisting}[language=Rust, caption={transaction.rs (service) --- The mutex and
map are one unit; the guard releases on drop.}]
let mut map = per_type.lock().expect("poisoned mutex");
*map.entry(*tx.transaction_type()).or_insert(0) += value;
\end{lstlisting}

For the map \emph{data structure} itself, C has no standard dictionary. Gentlix implements
\texttt{DayTotalMap} as a growable array with linear search; Rust uses
\texttt{BTreeMap<NaiveDate, i64>} from the standard library:

\begin{lstlisting}[language=C, caption={transaction\_service.c --- Ad hoc per-day map:
linear scan, then merge under \texttt{per\_day\_mutex}.}]
typedef struct DayTotalMap {
    DayTotal *entries;
    size_t    count;
    size_t    capacity;
} DayTotalMap;

static bool day_total_map_add(DayTotalMap *m,
                              int32_t date_key, int64_t amount) {
    for (size_t i = 0; i < m->count; ++i) {
        if (m->entries[i].date_key == date_key) {
            m->entries[i].total += amount;
            return true;
        }
    }
    /* grow array and append new date_key ... */
}
\end{lstlisting}

\begin{lstlisting}[language=Rust, caption={transaction.rs (service) --- Standard
map; mutex wraps updates in the worker loop.}]
let per_day = Arc::new(Mutex::new(BTreeMap::<NaiveDate, i64>::new()));
/* inside worker */
let mut map = per_day.lock().expect("poisoned mutex");
*map.entry(recorded.date_naive()).or_insert(0) += net;
\end{lstlisting}

After thread-local accumulation, C merges per-day totals under the same mutex pattern
as per-type totals:

\begin{lstlisting}[language=C, caption={transaction\_service.c --- Per-day merge
after the hot loop.}]
pthread_mutex_lock(&s->per_day_mutex);
for (size_t d = 0; d < local_day.count; ++d) {
    day_total_map_add(&s->per_day,
                      local_day.entries[d].date_key,
                      local_day.entries[d].total);
}
pthread_mutex_unlock(&s->per_day_mutex);
\end{lstlisting}

Schreiner~\cite{Schreiner1993} would call the C structure a hand-rolled collection;
the concurrency pattern is identical: local map first, shared map under mutex second.

\subsection{Atomics and counter totals}
\label{subsec:ch5sec3atomics}

Global counters (\texttt{total\_incoming}, \texttt{total\_outgoing}, volume, transaction
count) are updated from every worker simultaneously. C uses \texttt{\_Atomic int64\_t}
with \texttt{memory\_order\_relaxed}; Rust uses
\texttt{AtomicI64::fetch\_add(..., Ordering::Relaxed)}. Bos~\cite{Bos2023} explains that
relaxed ordering suffices when each atomic accumulates independently and no other memory
read depends on ordering---display totals in Gentlix, not ledger commits.

The hot loop in C increments shared counters with relaxed atomics:

\begin{lstlisting}[language=C, caption={transaction\_service.c --- Relaxed atomics in
the per-transaction loop.}]
atomic_fetch_add_explicit(&s->transaction_count, 1, memory_order_relaxed);
if (net > 0) {
    atomic_fetch_add_explicit(&s->total_incoming, net,
                              memory_order_relaxed);
} else if (net < 0) {
    atomic_fetch_add_explicit(&s->total_outgoing, -net,
                              memory_order_relaxed);
}
\end{lstlisting}

Rust mirrors the same arithmetic on \texttt{AtomicI64}:

\begin{lstlisting}[language=Rust, caption={transaction.rs (service) --- Same counters,
same ordering.}]
transaction_count.fetch_add(1, Ordering::Relaxed);
if net > 0 {
    total_incoming.fetch_add(net, Ordering::Relaxed);
} else if net < 0 {
    total_outgoing.fetch_add(-net, Ordering::Relaxed);
}
\end{lstlisting}

Plain \texttt{int64\_t} shared across C threads without atomics would be a data race with
undefined behaviour. Rust rejects unsynchronised \texttt{\&mut} sharing unless the type
is \texttt{Sync}. McNamara~\cite{McNamara2021} contrasts this with manually reasoning about
every load and store in C---the same numeric result when correct, a different proof burden.

\subsection{Threads, scopes, and join discipline}
\label{subsec:ch5sec3threads}

Workers borrow slices of a parent-owned transaction array. C passes raw pointers into
\texttt{pthread\_create}; the parent must \texttt{pthread\_join} every worker before
\texttt{free\_tx\_array}:

\begin{lstlisting}[language=C, caption={transaction\_service.c --- Workers must finish
before the shared array is freed.}]
for (size_t t = 0; t < num_threads; ++t)
    pthread_join(threads[t], NULL);
free_tx_array(all_txs, all_count);
\end{lstlisting}

If \texttt{pthread\_create} fails halfway, Gentlix joins started threads, destroys mutexes,
frees arrays, and returns a concurrency error:

\begin{lstlisting}[language=C, caption={transaction\_service.c --- Partial spawn
cleanup before returning \texttt{SERVICE\_ERROR\_CONCURRENCY}.}]
if (pthread_create(&threads[t], NULL, stats_worker_fn, &args[t]) != 0) {
    for (size_t j = 0; j < t; ++j) pthread_join(threads[j], NULL);
    free(args); free(threads); free_tx_array(all_txs, all_count);
    if (err) *err = service_error_internal(
        "failed to spawn statistics worker thread");
    return false;
}
\end{lstlisting}

\begin{lstlisting}[language=Rust, caption={transaction.rs (service) --- Join failure
maps to \texttt{ServiceError::Concurrency}.}]
stats_handle.await.map_err(|e|
    ServiceError::Concurrency(format!("join error: {e}")))?;
\end{lstlisting}

Blackburn~\cite{Blackburn2025}
classifies ``worker outlives buffer'' as temporal unsafe behaviour that can yield wrong
dashboard totals under load.

Rust uses \texttt{thread::scope} so workers cannot outlive the borrowed slice; join order
is enforced by the API:

\begin{lstlisting}[language=Rust, caption={transaction.rs (service) ---
\texttt{thread::scope} ties worker lifetimes to the transaction chunk.}]
thread::scope(|scope| {
    for chunk in all_transactions.chunks(chunk_size) {
        scope.spawn(move || {
            for tx in chunk { /* filters, atomics, mutex maps */ }
        });
    }
});
\end{lstlisting}

\subsection{The analytics pipeline end-to-end}
\label{subsec:ch5sec3pipeline}

The three mechanisms above appear together in one flow. After loading rows per owned
account, both backends deduplicate internal transfers before spawning workers. C sorts
pointers by id and compacts duplicates in place:

\begin{lstlisting}[language=C, caption={transaction\_service.c --- Dedupe runs before
any \texttt{pthread\_create} call.}]
static size_t dedupe_transactions_by_id(Transaction **all_txs, size_t all_count) {
    if (all_count <= 1) return all_count;
    qsort(all_txs, all_count, sizeof(Transaction *), compare_tx_ptr_by_id);
    size_t w = 1;
    for (size_t i = 1; i < all_count; ++i) {
        if (transaction_id(all_txs[i]).value ==
            transaction_id(all_txs[w - 1]).value) {
            transaction_free(all_txs[i]);   // drop duplicate
        } else {
            all_txs[w++] = all_txs[i];
        }
    }
    return w;
}
\end{lstlisting}

Rust fills a \texttt{HashMap<TransactionId, Transaction>} while iterating accounts:

\begin{lstlisting}[language=Rust, caption={transaction.rs (service) --- Last write wins
for duplicate ids across account lists.}]
let mut unique: HashMap<TransactionId, Transaction> = HashMap::new();
for account in &accounts {
    let txs = self.tx_repo.list_for_account(account.id().unwrap(),
        per_account_limit, 0).await?;
    for tx in txs {
        if let Some(id) = tx.id() {
            unique.entry(id).or_insert(tx);
        }
    }
}
\end{lstlisting}

Dedupe stays single-threaded so parallel code never double-counts a row seen from two
accounts. Each worker then applies the same date, type, and ownership filters before
updating counters:

\begin{lstlisting}[language=C, caption={transaction\_service.c --- Filter logic in
\texttt{stats\_worker\_fn} hot loop.}]
time_t rec = transaction_recorded_on(tx);
if (arg->has_from && rec < arg->from) continue;
if (arg->has_to   && rec > arg->to)   continue;
if (arg->has_tx_type &&
    transaction_type_get(tx) != arg->filter_type) continue;
bool to_owned = id_in_set(arg->user_ids, arg->user_count,
                          transaction_to_account_id(tx).value);
bool from_owned = id_in_set(arg->user_ids, arg->user_count,
                            transaction_from_account_id(tx).value);
int64_t net = (to_owned ? value : 0) - (from_owned ? value : 0);
\end{lstlisting}

\begin{lstlisting}[language=Rust, caption={transaction.rs (service) --- Same filters
inside the scoped worker.}]
if let Some(f) = from { if recorded < f { continue; } }
if let Some(t) = to   { if recorded > t { continue; } }
if let Some(filter_type) = filters.transaction_type {
    if *tx.transaction_type() != filter_type { continue; }
}
let net = (if to_owned { value } else { 0 })
        - (if from_owned { value } else { 0 });
\end{lstlisting}

Palmieri~\cite{Palmieri2022} would classify the map key as a newtype
(\texttt{TransactionId}) that prevents confusing ids with raw integers---the same id
type used in Section~\ref{sec:ch5sec1}.

Worker count follows \texttt{get\_processor\_count()} in C and
\texttt{thread::available\_parallelism()} in Rust:

\begin{lstlisting}[language=C, caption={transaction\_service.c --- Thread count from
the host OS.}]
static long get_processor_count(void) {
#ifdef _WIN32
    SYSTEM_INFO info;
    GetSystemInfo(&info);
    return (long)info.dwNumberOfProcessors;
#else
    long n = sysconf(_SC_NPROCESSORS_ONLN);
    return n > 0 ? n : 1;
#endif
}
\end{lstlisting}

\begin{lstlisting}[language=Rust, caption={transaction.rs (service) --- Same intent.}]
let num_threads = thread::available_parallelism()
    .map(|n| n.get())
    .unwrap_or(1);
let chunk_size = (all_transactions.len() / num_threads).max(1);
\end{lstlisting}

Each worker applies date, type, and ownership filters, updates relaxed atomics, fills
local maps, then merges under the mutexes from Section~\ref{subsec:ch5sec3mutex}.
Rust wraps the entire CPU-heavy phase in \texttt{task::spawn\_blocking}; C runs analytics
inline in the HTTP handler before the response is sent:

\begin{lstlisting}[language=C, caption={transaction\_service.c --- No runtime offload;
\texttt{pthread} workers run in the request thread's call stack.}]
bool transaction_service_compute_user_statistics(
    TransactionService *svc, UserId user_id,
    int64_t per_account_limit, /* ... */, ServiceError *err)
{
    /* load accounts, dedupe, pthread_create / join, aggregate */
}
\end{lstlisting}

\begin{lstlisting}[language=Rust, caption={transaction.rs (service) --- CPU work
offloaded from the Tokio worker.}]
let stats_handle = task::spawn_blocking(move || {
    thread::scope(|scope| { /* workers */ });
    ExtendedUserStatistics { /* ... */ }
});
stats_handle.await.map_err(|e|
    ServiceError::Concurrency(format!("join error: {e}")))?;
\end{lstlisting}

Flitton and Morton~\cite{Flitton2023}
warn that long CPU sections on a Tokio worker stall other async tasks---offloading is a
runtime concern (Section~\ref{subsec:ch6sec2-4}), while mutex/atomic/thread rules in this
section describe the algorithm both languages share.

\section*{Summary}

The three sections of this chapter show that Gentlix enforces banking correctness through
types, explicit error chains, and a deliberate parallel aggregation design. In C, every
invariant and every lock is a convention: validated email can still look like a raw
string, \texttt{ServiceError *err} can be ignored, and mutex init/destroy must be remembered
on each path. Rust moves much of that obligation to compile time---newtypes, \texttt{Result},
\texttt{thread::scope}, and \texttt{Mutex<T>}---while implementing the same product rules.

Database column indices and HTTP JSON envelopes are not repeated here; they belong to
Chapter~\ref{sec:ch6sec1} and Section~\ref{subsec:ch6sec2-5}. The service layer assumes
correct persistence and a faithful HTTP mapping; this chapter shows how the two backends
justify that trust before a request leaves \texttt{service/}.
